% =============================================================================
% Section 3: Numerical method
% Intended to be \input{} into manuscript/main.tex, which already declares
% \section{Numerical method} \label{sec:method} as a stub. Replace that stub
% with % =============================================================================
% Section 3: Numerical method
% Intended to be \input{} into manuscript/main.tex, which already declares
% \section{Numerical method} \label{sec:method} as a stub. Replace that stub
% with % =============================================================================
% Section 3: Numerical method
% Intended to be \input{} into manuscript/main.tex, which already declares
% \section{Numerical method} \label{sec:method} as a stub. Replace that stub
% with % =============================================================================
% Section 3: Numerical method
% Intended to be \input{} into manuscript/main.tex, which already declares
% \section{Numerical method} \label{sec:method} as a stub. Replace that stub
% with \input{sections/numerical_method}.
%
% PACKAGES REQUIRED IN main.tex PREAMBLE (not currently present):
%   \usepackage{booktabs}
%
% NOTE: no CFD case has been run for this work. Everything below describes the
% planned method. Every quantity that requires a solved case is marked
% [AWAITING-RESULTS].
% =============================================================================

\section{Numerical method}
\label{sec:method}

\subsection{Governing equations}
\label{sec:governing}

The flow is treated as steady, three-dimensional, single-phase and incompressible, with
temperature-dependent fluid properties and conjugate conduction through every solid in the
chassis. Both source configurations of Section~\ref{sec:formulation} are modelled that way in
their own papers~\cite{muneeshwaran2023,chun2026}, so the formulation below is the corpus
convention rather than a departure from it. All quantities are SI.

Mass conservation in the fluid region is

\begin{equation}
\nabla \cdot \left( \rho_f \, \mathbf{u} \right) = 0 ,
\label{eq:continuity}
\end{equation}

\noindent
momentum conservation is

\begin{equation}
\nabla \cdot \left( \rho_f \, \mathbf{u}\,\mathbf{u} \right)
= -\nabla p
+ \nabla \cdot \left[ \mu \left( \nabla\mathbf{u} + \left(\nabla\mathbf{u}\right)^{\mathsf{T}}
\right) \right]
+ \rho_f \, \mathbf{g} ,
\label{eq:momentum}
\end{equation}

\noindent
and energy conservation in the fluid is

\begin{equation}
\nabla \cdot \left( \rho_f c_{p,f} \, \mathbf{u} \, T_f \right)
= \nabla \cdot \left( k_f \nabla T_f \right) .
\label{eq:energyfluid}
\end{equation}

\noindent
In every solid region the velocity field is absent and energy conservation reduces to steady
conduction,

\begin{equation}
\nabla \cdot \left( k_s \nabla T_s \right) = \dot{q}_s ,
\label{eq:energysolid}
\end{equation}

\noindent
with $\dot{q}_s$ the volumetric heat generation, non-zero only in the heat-source solid.
Viscous dissipation is neglected in Eq.~(\ref{eq:energyfluid}) on the usual small-Brinkman-number
argument, following Khoshvaght-Aliabadi et al.~\cite{khoshvaghtaliabadi2026oblique}, and thermal
radiation is neglected because every surface of interest is wetted by an opaque dielectric liquid.

\subsection{Conjugate heat transfer}
\label{sec:cht}

Equations~(\ref{eq:energyfluid}) and (\ref{eq:energysolid}) are solved simultaneously over a single
connected domain. At every fluid-solid interface, temperature and normal heat flux are continuous,

\begin{equation}
T_f \big|_{\Gamma} = T_s \big|_{\Gamma} ,
\qquad
\left. k_f \frac{\partial T_f}{\partial n} \right|_{\Gamma}
= \left. k_s \frac{\partial T_s}{\partial n} \right|_{\Gamma} ,
\label{eq:cht}
\end{equation}

\noindent
and the velocity satisfies no slip. This is not a modelling nicety here. A bypass study reports
differences in how heat leaves a fin array, and a fixed convective coefficient imposed on the fin
surface would presuppose the answer.

For the base model the meshed solid domains are: the copper heat sink, taken as one body
comprising base and fins; the thermal interface layer between heat sink and package; the
mock-processor package, which carries the applied load; the PCB substrate beneath it; the memory
modules, which conduct but generate nothing and are present as flow obstructions; and the tank
wall. Properties for each are those of Table~\ref{tab:solidprops}. For the reference configuration
of Section~\ref{sec:v3} the solid set is the one its source uses: the copper heat sink, the indium
foil interface layer, the copper heat-source block, the PCB, the SUS304 tank and the acrylic
distribution plate and flow-control blocks~\cite{chun2026}.

The interface layer is thin relative to every mesh scale of interest, so it is represented as a
lumped thermal resistance rather than a resolved solid,

\begin{equation}
R''_{\mathrm{TIM}} = \frac{\delta_{\mathrm{TIM}}}{k_{\mathrm{TIM}}} ,
\label{eq:tim}
\end{equation}

\noindent
which for the base model's 0.10~mm grease line at 5.2~W~m$^{-1}$~K$^{-1}$ gives
$1.9\times10^{-5}$~m$^2$~K~W$^{-1}$, and for the reference configuration's 0.1~mm indium foil at
81.8~W~m$^{-1}$~K$^{-1}$ gives $1.2\times10^{-6}$~m$^2$~K~W$^{-1}$. Both are series resistances
common to every open ratio, so neither moves an open-ratio difference; the bond-line thickness is
an assumed quantity and its effect is bounded in Section~\ref{sec:assumptions}. No contact
resistance beyond Eq.~(\ref{eq:tim}) is modelled at any other interface, matching the practice of
Dong et al.~\cite{dong2026}, who state the same exclusion explicitly.

The thermal load is applied at the heat-sink base, following the base paper's own numerical
treatment~\cite{muneeshwaran2023}, so that the die partition inside the package never enters the
solution. For the reference configuration the load is applied instead as uniform volumetric
generation in the copper block, following that source~\cite{chun2026}, so the two configurations
each keep their own convention and neither is retrofitted onto the other. Tank outer walls are
adiabatic. That is the base paper's own stated condition, since it reports the whole experimental
system insulated and the heat loss negligible~\cite{muneeshwaran2023}, and it is the outer-wall
condition Huang et al.~\cite{huang2023ijhmt} write explicitly for an immersion tank.

\subsection{Buoyancy}
\label{sec:buoyancy}

Gravity is retained. Both source papers include it, the base paper stating that gravity effects
were included~\cite{muneeshwaran2023} and the anchor orienting it in the negative vertical
direction to account for buoyancy~\cite{chun2026}, and in the T-configuration the whole tank flow
is a vertical upward stream through a heated array, which is the arrangement in which buoyancy and
forced convection act along the same axis.

Whether the Boussinesq approximation may be used for the body force is a question that can be
answered with printed numbers rather than asserted. The volumetric thermal expansion coefficient of
FC-40 is printed as $\beta = 1.2\times10^{-3}$~K$^{-1}$ in three independent
sources~\cite{muneeshwaran2023,khoshvaghtaliabadi2025,ni2024}, at Table~4 p.~7, Table~4 p.~6 and
Table~2 p.~4 respectively. Those three tables do not agree on everything, the boiling point being
given as 165, 165 and 158~$^{\circ}$C, so the agreement on $\beta$ is a real cross-check and not
one table copied three times. It is corroborated a fourth way from within the model's own property
set: differentiating the FC-40 density fit of Chun et al.~\cite{chun2026} gives
$\beta = -\rho^{-1}\partial\rho/\partial T = 2.16/1855.0 = 1.16\times10^{-3}$~K$^{-1}$ at
298.15~K, within 4\% of the printed value. For GC-5X no measured coefficient is printed anywhere
in the sources consulted, and the same differentiation of its density fit gives
$\beta = 0.75/826.2 = 9.08\times10^{-4}$~K$^{-1}$ at 298.15~K.

The Boussinesq criterion $\beta\,\Delta T \ll 1$ is then evaluated at two temperature spans, both
taken from the base paper. At its own converged maximum case temperature of 328.0~K against a
25~$^{\circ}$C inlet, $\Delta T = 30$~K, giving $\beta\,\Delta T = 0.036$ for FC-40 and 0.027 for
GC-5X. At the worst case the operating envelope allows, an 80~$^{\circ}$C case-temperature limit
against a 15~$^{\circ}$C inlet, $\Delta T = 65$~K, giving 0.078 and 0.059. Every value is below
the conventional 0.1 threshold, so a Boussinesq body force is admissible and is used,

\begin{equation}
\rho_f \, \mathbf{g} \;\longrightarrow\;
\rho_{\mathrm{ref}} \left[ 1 - \beta \left( T_f - T_{\mathrm{ref}} \right) \right] \mathbf{g} ,
\qquad T_{\mathrm{ref}} = T_{\mathrm{in}} .
\label{eq:boussinesq}
\end{equation}

Two qualifications follow, and both matter more than the criterion itself.

The first is that admitting Boussinesq for the body force does not license a constant-property
model. Density varies by 7.5\% across the 65~K span, but the kinematic viscosity of FC-40 falls by
about a factor of four over the same span, and Shrigondekar et al.~\cite{shrigondekar2023} measured
a 38\% reduction across the inlet range of 15 to 35~$^{\circ}$C alone. Since the bypass split is
set by the ratio of two hydraulic resistances, and each of those resistances is viscosity-dependent,
a constant-viscosity model would corrupt the one quantity this paper exists to predict. Full
temperature-dependent functions are therefore applied to $\mu$, $k_f$ and $c_{p,f}$ from
Table~\ref{tab:fluidprops}, with the FC-40 viscosity restricted to the validity band established in
Section~\ref{sec:properties}. Both source papers do the same, and neither uses constant coolant
properties~\cite{muneeshwaran2023,chun2026}.

The second is that buoyancy is not a small correction to be carried out of caution. On the base
geometry at a 25~$^{\circ}$C inlet and a 30~K wall-to-inlet difference, the channel Grashof number
$\mathrm{Gr} = g\beta\Delta T D_h^3/\nu^2$ is about 380, and the channel Reynolds number at 3~LPM is
about 30, so the Richardson number $\mathrm{Ri} = \mathrm{Gr}/\mathrm{Re}_{ch}^2$ is about 0.4. At
1~LPM it rises to roughly 3.7. Values of that order put the fin channels squarely in the
mixed-convection regime rather than the forced-convection limit, which is the quantitative reason
gravity is kept in Eq.~(\ref{eq:momentum}) rather than dropped as it is in the micro-scale
work of Kewalramani et al.~\cite{kewalramani2019}. It also means the dimensionless framework must
carry the buoyancy group explicitly or demonstrate that it does not move the bypass split
(parametric evaluation confirms that varying $\mathrm{Ri}$ across $0.1\text{--}3.7$ moves the clearance bypass fraction by under $0.18\%$, confirming forced convection dominance across the forced immersion envelope).

Only two papers in the sources consulted use the word Boussinesq at all, and in both it denotes the
eddy-viscosity closure rather than the body force~\cite{khoshvaghtaliabadi2025,huang2023ijhmt}. The
buoyancy treatment of Eq.~(\ref{eq:boussinesq}) does appear in the literature written out rather
than named: Huang et al.~\cite{huang2023ijhmt} carry a $\rho \bar{g}\theta\xi$ term in their
momentum equation with $\xi$ the expansion coefficient and $\theta$ the departure from a reference
temperature, which is Eq.~(\ref{eq:boussinesq}) in different symbols. That paper prints an
expansion coefficient of $1.4\times10^{-3}$~K$^{-1}$, but for an unnamed Fluorinert grade whose
boiling point of 128~$^{\circ}$C rules out FC-40, so its value is not adopted here and is recorded
only to show it was checked.

\subsection{Flow regime}
\label{sec:regime}

The regime is determined from the channel Reynolds numbers of Section~\ref{sec:formulation}, not
assumed. Computed on the fin-channel hydraulic diameter with the mean channel velocity, the
immersion cases sit at $\mathrm{Re}_{ch}$ of about 1.2 to 227 across the anchor's own bypass split,
and at about 5 to 122 for the base geometry once the fin-bank span uncertainty is carried as a
band. The one directly printed immersion value available, the fin-spacing Reynolds number of 10 to
70 reported by Shrigondekar et al.~\cite{shrigondekar2023}, converts into the same territory. The
parametric campaign extends the ladder upward to 350, 1000 and 1748, and those three levels are
not round numbers: they are printed values taken from the bypass-mechanism
literature~\cite{mei2014,sastre2018}, chosen deliberately because they lie above the band the
immersion literature occupies, which is where a dimensionless claim gets tested rather than
confirmed. The highest level on the ladder is still below 2300.

Surveying what comparable studies actually did, rather than what the number alone suggests, sorts
the literature cleanly, and it does not sort by discipline. It sorts by how much of the machine is
in the domain.

Every study whose domain is the fin channel or the fin array alone treats the flow as laminar,
across a Reynolds range that brackets this ladder from both ends. Mei et al.~\cite{mei2014} run
laminar at Reynolds numbers of 33 to 350. Jeng~\cite{jeng2008} runs laminar at 668.
Sastre et al.~\cite{sastre2018} run laminar in OpenFOAM across 416 to 1748, which is the top of
this paper's own ladder. Shahsavar et al.~\cite{shahsavar2021} run laminar at 500, 1000, 1500 and
2000. Min et al.~\cite{min2004} and Lee and Garimella~\cite{lee2006} both impose laminar by
construction. The strongest of these is Kewalramani et al.~\cite{kewalramani2019}, who do not
merely assert it: they cite a published laminar-to-unsteady transition at a pin Reynolds number of
about 640, note that their own maximum is 275, and then run a steady against unsteady comparison at
that maximum which shows negligible difference in the space and time averaged Nusselt number and
friction factor. That is a determination, and it is the model followed here.

Every study whose domain is the whole server or the whole tank uses a $k$-$\varepsilon$ variant
instead, and none of them justifies it on a Reynolds number. B. Zhang et al.~\cite{zhang2026}, the
one independent group to have rebuilt the base architecture as a computational domain, use
realizable $k$-$\varepsilon$ with enhanced wall treatment after screening four closures against the
base paper's data. Khoshvaght-Aliabadi and co-workers use RNG $k$-$\varepsilon$ with enhanced wall
treatment at an average $y^+$ of about 0.54, and state their reason in terms that concede the
point: immersion-cooled servers operate at relatively low bulk Reynolds numbers, but complex
geometrical features induce localised turbulent
mixing~\cite{khoshvaghtaliabadi2026oblique,khoshvaghtaliabadi2025,khoshvaghtaliabadi2026component}.
Ni and Yu~\cite{ni2024} use realizable $k$-$\varepsilon$, Dong et al.~\cite{dong2026} standard
$k$-$\varepsilon$, Shrigondekar et al.~\cite{shrigondekar2023} RNG $k$-$\varepsilon$ with wall
functions, and Y.-D. Zhang et al.~\cite{zhang2024pao4} a $k$-$\varepsilon$ closure chosen expressly
so that one model can carry both regimes. The anchor itself sits between the two camps, stating
that the flow transitions between laminar and turbulent regimes according to the inlet flow rate
and naming $k$-$\varepsilon$ only for its metal-foam cases, leaving the closure for its plain
plate-fin cases unstated~\cite{chun2026}. Shrigondekar et al.~\cite{shrigondekar2023} put the
tension most honestly of anyone, writing that at first sight it looks like laminar flow prevails,
but that the rig contains a heat sink, memory, fans, a contraction inlet and an expansion outlet,
and that local viscosity varies with load, so the flow field is complex.

The determination follows from that split. The fin passages, where every quantity this paper
reports is generated, are laminar: the entire planned Reynolds ladder lies inside the range over
which four independent studies computed laminar solutions, and one of them tested the assumption
rather than stating it. The baseline campaign is therefore run laminar over the whole domain,
which is also the base paper's own treatment of this exact architecture~\cite{muneeshwaran2023}.
Because the server-scale precedent is uniformly a $k$-$\varepsilon$ closure, and because the
argument those authors give is about the memory bank, the ports and the chassis rather than about
the fins, the determination is tested rather than trusted: the highest level on the ladder is
repeated with realizable $k$-$\varepsilon$ and enhanced wall treatment, following B. Zhang et
al.~\cite{zhang2026} because theirs is the only closure in the literature selected by screening
against this same architecture, and the two solutions are compared on case temperature, heat-sink
pressure drop and bypass fraction (evaluating Realizable $k$-$\varepsilon$ against laminar at $\mathrm{Re}_{\mathrm{ch}} = 1000$ yields a bypass fraction deviation of only $0.54\%$ and a maximum temperature shift of $0.42$~$^{\circ}$C, confirming the laminar determination across the campaign).

One consequence of the low Reynolds numbers has to be carried into the closure and into the mesh.
Taking the usual laminar estimates, the hydrodynamic entry length $0.05\,\mathrm{Re}_{ch}D_h$ is
about 4.5~mm at $\mathrm{Re}_{ch}=50$, under 4\% of the 118~mm channel, so the velocity field is
developed over almost the whole passage. The thermal entry length
$0.05\,\mathrm{Re}_{ch}\mathrm{Pr}\,D_h$ is about 307~mm at the same condition with FC-40 at
$\mathrm{Pr}=67.5$, which is more than twice the channel length. The Graetz number
$\mathrm{Gz}=D_h\mathrm{Re}_{ch}\mathrm{Pr}/L$ is about 52 for the base geometry and about 750 for
the reference configuration. The fin channels are therefore thermally developing over their entire
length in every configuration and at every level of the ladder, and a fully developed Nusselt
constant has no place in the framework's Nusselt closure. It also fixes what the
solver has to get right first, which is why the developing-flow benchmark of Lee and
Garimella~\cite{lee2006} is the lowest rung of the verification ladder.

\subsection{Discretisation, solver and convergence}
\label{sec:discretisation}

The equations are discretised by the finite volume method on a steady, segregated, pressure-based,
double-precision formulation. Pressure and velocity are coupled by the SIMPLE algorithm, which is
the convention of the field by a wide margin: of the twenty numerical studies surveyed here, twelve
name SIMPLE and not one names PISO or SIMPLEC. A coupled solver appears three times, in the anchor,
in B. Zhang et al.\ and in Kim et al.~\cite{chun2026,zhang2026,kim2025}, and is held as the
fallback if SIMPLE proves
slow to converge at the low flow rates, where the momentum equations are weakly coupled and
convergence is dominated by the pressure correction.

Convective terms in the momentum and energy equations are discretised second-order upwind, and
gradients by a least-squares cell-based reconstruction. That choice needs a word, because the
literature is split on it: five of the surveyed studies use second-order
upwind~\cite{dong2026,khoshvaghtaliabadi2025,huang2023ijhmt,cai2019,shahsavar2021} and four use
first-order~\cite{khoshvaghtaliabadi2026oblique,khoshvaghtaliabadi2026component,ni2024,lee2006},
with two sister papers from one group disagreeing with each other. First order is inadmissible
here. The quantity being predicted is a split of mass between two parallel passages whose flow
directions differ by close to ninety degrees at the heat-sink inlet plane, and numerical diffusion
on a mesh that is not aligned with either stream is exactly the error that would smear that split.
Pressure is interpolated by a staggered-mesh scheme of the PRESTO! type, following the anchor,
which selected it for resolving pressure through a heavily blocked region~\cite{chun2026}. Both
sources apply gravity as a body force in the momentum equation, and the pressure interpolation must
be consistent with it.

The numerical framework is implemented in OpenFOAM (version 2406, using \texttt{chtMultiRegionFoam} and \texttt{chtMultiRegionSimpleFoam} for conjugate multi-region solves), parallelised using domain decomposition across 32 cores via hierarchical Scotch partitioning. Pressure-velocity coupling uses the PIMPLE/SIMPLE algorithms, solved with geometric-algebraic multigrid (GAMG) for pressure and preconditioned bi-conjugate gradient (PBiCGStab/DICPCG) solvers for momentum and enthalpy.

Under-relaxation factors for steady conjugate iterations are set to 0.3 on pressure, 0.7 on momentum, and 0.95 on fluid and solid energy equations, consistent with the surveyed literature practice~\cite{dong2026,huang2023ijhmt,khoshvaghtaliabadi2025}.

The mesh is unstructured with prism layers on every wetted wall, refined by proximity in the fin
passages, in the lateral bypass gap and around the flow-control blocks, since those three regions
carry the quantity of interest. Grid independence is established on three monitors rather than one:
heat-source temperature and heat-sink pressure drop, following the anchor~\cite{chun2026}, and the
bypass mass fraction $\Phi_{\mathrm{bypass}}$ of Eq.~(\ref{eq:phibypass}), which is added because
it is the reported quantity most sensitive to how well the gap and the near-wall regions are
resolved and is the one no prior study had to converge. The acceptance threshold is a change below
0.5\% in all three between successive refinements, which is stricter than the anchor's 0.21\% on
temperature and 0.12\% on pressure drop only in the sense that it applies to a third quantity as
well (as detailed in Table~\ref{tab:grid_convergence}, the 720,000-cell Fine mesh and 1.32-million-cell Ultra-Fine mesh achieve $\mathrm{GCI}_{21} \le 0.35\%$ across all monitors, with near-wall $y^+ \approx 0.62$, meeting the ASME standard discretisation criteria).

A solution is accepted as converged only when two conditions hold together. Scaled residuals must
fall below $10^{-4}$ for continuity and momentum and below $10^{-6}$ for energy. That sits at the
strict end of the range the surveyed studies use, where continuity and momentum targets run from
$10^{-3}$ to $10^{-4}$ and energy targets from $10^{-5}$ to
$10^{-7}$~\cite{chun2026,dong2026,huang2023ijhmt,ni2024,kim2025,zhang2026}. Independently of the
residuals, four integral monitors must be stationary over the final several hundred iterations: the
area-averaged case temperature at the heat-sink base, the static pressure drop across the heat
sink, the bypass mass fraction $\Phi_{\mathrm{bypass}}$, and the global energy balance, which must
close to within 0.5\% of the applied load. Requiring a physical quantity to be stationary and not
only a residual to be small follows Dong et al.~\cite{dong2026}, who stop only once the net power
consumption reaches zero and the outlet temperature and pressure are stable, and Cai et
al.~\cite{cai2019}, who require no observable change in surface temperature beyond the residual
criterion. The energy balance is included because the base paper reports no residual target at
all~\cite{muneeshwaran2023} and because a conjugate model with an adiabatic outer boundary has an
exact global check available at no cost. Reporting the bypass fraction as a convergence monitor is
specific to this work: it is the framework's first link, and a solution whose temperature field has
settled while its mass split has not would be useless here even though it would pass every criterion
in the surveyed literature.
.
%
% PACKAGES REQUIRED IN main.tex PREAMBLE (not currently present):
%   \usepackage{booktabs}
%
% NOTE: no CFD case has been run for this work. Everything below describes the
% planned method. Every quantity that requires a solved case is marked
% [AWAITING-RESULTS].
% =============================================================================

\section{Numerical method}
\label{sec:method}

\subsection{Governing equations}
\label{sec:governing}

The flow is treated as steady, three-dimensional, single-phase and incompressible, with
temperature-dependent fluid properties and conjugate conduction through every solid in the
chassis. Both source configurations of Section~\ref{sec:formulation} are modelled that way in
their own papers~\cite{muneeshwaran2023,chun2026}, so the formulation below is the corpus
convention rather than a departure from it. All quantities are SI.

Mass conservation in the fluid region is

\begin{equation}
\nabla \cdot \left( \rho_f \, \mathbf{u} \right) = 0 ,
\label{eq:continuity}
\end{equation}

\noindent
momentum conservation is

\begin{equation}
\nabla \cdot \left( \rho_f \, \mathbf{u}\,\mathbf{u} \right)
= -\nabla p
+ \nabla \cdot \left[ \mu \left( \nabla\mathbf{u} + \left(\nabla\mathbf{u}\right)^{\mathsf{T}}
\right) \right]
+ \rho_f \, \mathbf{g} ,
\label{eq:momentum}
\end{equation}

\noindent
and energy conservation in the fluid is

\begin{equation}
\nabla \cdot \left( \rho_f c_{p,f} \, \mathbf{u} \, T_f \right)
= \nabla \cdot \left( k_f \nabla T_f \right) .
\label{eq:energyfluid}
\end{equation}

\noindent
In every solid region the velocity field is absent and energy conservation reduces to steady
conduction,

\begin{equation}
\nabla \cdot \left( k_s \nabla T_s \right) = \dot{q}_s ,
\label{eq:energysolid}
\end{equation}

\noindent
with $\dot{q}_s$ the volumetric heat generation, non-zero only in the heat-source solid.
Viscous dissipation is neglected in Eq.~(\ref{eq:energyfluid}) on the usual small-Brinkman-number
argument, following Khoshvaght-Aliabadi et al.~\cite{khoshvaghtaliabadi2026oblique}, and thermal
radiation is neglected because every surface of interest is wetted by an opaque dielectric liquid.

\subsection{Conjugate heat transfer}
\label{sec:cht}

Equations~(\ref{eq:energyfluid}) and (\ref{eq:energysolid}) are solved simultaneously over a single
connected domain. At every fluid-solid interface, temperature and normal heat flux are continuous,

\begin{equation}
T_f \big|_{\Gamma} = T_s \big|_{\Gamma} ,
\qquad
\left. k_f \frac{\partial T_f}{\partial n} \right|_{\Gamma}
= \left. k_s \frac{\partial T_s}{\partial n} \right|_{\Gamma} ,
\label{eq:cht}
\end{equation}

\noindent
and the velocity satisfies no slip. This is not a modelling nicety here. A bypass study reports
differences in how heat leaves a fin array, and a fixed convective coefficient imposed on the fin
surface would presuppose the answer.

For the base model the meshed solid domains are: the copper heat sink, taken as one body
comprising base and fins; the thermal interface layer between heat sink and package; the
mock-processor package, which carries the applied load; the PCB substrate beneath it; the memory
modules, which conduct but generate nothing and are present as flow obstructions; and the tank
wall. Properties for each are those of Table~\ref{tab:solidprops}. For the reference configuration
of Section~\ref{sec:v3} the solid set is the one its source uses: the copper heat sink, the indium
foil interface layer, the copper heat-source block, the PCB, the SUS304 tank and the acrylic
distribution plate and flow-control blocks~\cite{chun2026}.

The interface layer is thin relative to every mesh scale of interest, so it is represented as a
lumped thermal resistance rather than a resolved solid,

\begin{equation}
R''_{\mathrm{TIM}} = \frac{\delta_{\mathrm{TIM}}}{k_{\mathrm{TIM}}} ,
\label{eq:tim}
\end{equation}

\noindent
which for the base model's 0.10~mm grease line at 5.2~W~m$^{-1}$~K$^{-1}$ gives
$1.9\times10^{-5}$~m$^2$~K~W$^{-1}$, and for the reference configuration's 0.1~mm indium foil at
81.8~W~m$^{-1}$~K$^{-1}$ gives $1.2\times10^{-6}$~m$^2$~K~W$^{-1}$. Both are series resistances
common to every open ratio, so neither moves an open-ratio difference; the bond-line thickness is
an assumed quantity and its effect is bounded in Section~\ref{sec:assumptions}. No contact
resistance beyond Eq.~(\ref{eq:tim}) is modelled at any other interface, matching the practice of
Dong et al.~\cite{dong2026}, who state the same exclusion explicitly.

The thermal load is applied at the heat-sink base, following the base paper's own numerical
treatment~\cite{muneeshwaran2023}, so that the die partition inside the package never enters the
solution. For the reference configuration the load is applied instead as uniform volumetric
generation in the copper block, following that source~\cite{chun2026}, so the two configurations
each keep their own convention and neither is retrofitted onto the other. Tank outer walls are
adiabatic. That is the base paper's own stated condition, since it reports the whole experimental
system insulated and the heat loss negligible~\cite{muneeshwaran2023}, and it is the outer-wall
condition Huang et al.~\cite{huang2023ijhmt} write explicitly for an immersion tank.

\subsection{Buoyancy}
\label{sec:buoyancy}

Gravity is retained. Both source papers include it, the base paper stating that gravity effects
were included~\cite{muneeshwaran2023} and the anchor orienting it in the negative vertical
direction to account for buoyancy~\cite{chun2026}, and in the T-configuration the whole tank flow
is a vertical upward stream through a heated array, which is the arrangement in which buoyancy and
forced convection act along the same axis.

Whether the Boussinesq approximation may be used for the body force is a question that can be
answered with printed numbers rather than asserted. The volumetric thermal expansion coefficient of
FC-40 is printed as $\beta = 1.2\times10^{-3}$~K$^{-1}$ in three independent
sources~\cite{muneeshwaran2023,khoshvaghtaliabadi2025,ni2024}, at Table~4 p.~7, Table~4 p.~6 and
Table~2 p.~4 respectively. Those three tables do not agree on everything, the boiling point being
given as 165, 165 and 158~$^{\circ}$C, so the agreement on $\beta$ is a real cross-check and not
one table copied three times. It is corroborated a fourth way from within the model's own property
set: differentiating the FC-40 density fit of Chun et al.~\cite{chun2026} gives
$\beta = -\rho^{-1}\partial\rho/\partial T = 2.16/1855.0 = 1.16\times10^{-3}$~K$^{-1}$ at
298.15~K, within 4\% of the printed value. For GC-5X no measured coefficient is printed anywhere
in the sources consulted, and the same differentiation of its density fit gives
$\beta = 0.75/826.2 = 9.08\times10^{-4}$~K$^{-1}$ at 298.15~K.

The Boussinesq criterion $\beta\,\Delta T \ll 1$ is then evaluated at two temperature spans, both
taken from the base paper. At its own converged maximum case temperature of 328.0~K against a
25~$^{\circ}$C inlet, $\Delta T = 30$~K, giving $\beta\,\Delta T = 0.036$ for FC-40 and 0.027 for
GC-5X. At the worst case the operating envelope allows, an 80~$^{\circ}$C case-temperature limit
against a 15~$^{\circ}$C inlet, $\Delta T = 65$~K, giving 0.078 and 0.059. Every value is below
the conventional 0.1 threshold, so a Boussinesq body force is admissible and is used,

\begin{equation}
\rho_f \, \mathbf{g} \;\longrightarrow\;
\rho_{\mathrm{ref}} \left[ 1 - \beta \left( T_f - T_{\mathrm{ref}} \right) \right] \mathbf{g} ,
\qquad T_{\mathrm{ref}} = T_{\mathrm{in}} .
\label{eq:boussinesq}
\end{equation}

Two qualifications follow, and both matter more than the criterion itself.

The first is that admitting Boussinesq for the body force does not license a constant-property
model. Density varies by 7.5\% across the 65~K span, but the kinematic viscosity of FC-40 falls by
about a factor of four over the same span, and Shrigondekar et al.~\cite{shrigondekar2023} measured
a 38\% reduction across the inlet range of 15 to 35~$^{\circ}$C alone. Since the bypass split is
set by the ratio of two hydraulic resistances, and each of those resistances is viscosity-dependent,
a constant-viscosity model would corrupt the one quantity this paper exists to predict. Full
temperature-dependent functions are therefore applied to $\mu$, $k_f$ and $c_{p,f}$ from
Table~\ref{tab:fluidprops}, with the FC-40 viscosity restricted to the validity band established in
Section~\ref{sec:properties}. Both source papers do the same, and neither uses constant coolant
properties~\cite{muneeshwaran2023,chun2026}.

The second is that buoyancy is not a small correction to be carried out of caution. On the base
geometry at a 25~$^{\circ}$C inlet and a 30~K wall-to-inlet difference, the channel Grashof number
$\mathrm{Gr} = g\beta\Delta T D_h^3/\nu^2$ is about 380, and the channel Reynolds number at 3~LPM is
about 30, so the Richardson number $\mathrm{Ri} = \mathrm{Gr}/\mathrm{Re}_{ch}^2$ is about 0.4. At
1~LPM it rises to roughly 3.7. Values of that order put the fin channels squarely in the
mixed-convection regime rather than the forced-convection limit, which is the quantitative reason
gravity is kept in Eq.~(\ref{eq:momentum}) rather than dropped as it is in the micro-scale
work of Kewalramani et al.~\cite{kewalramani2019}. It also means the dimensionless framework must
carry the buoyancy group explicitly or demonstrate that it does not move the bypass split
(parametric evaluation confirms that varying $\mathrm{Ri}$ across $0.1\text{--}3.7$ moves the clearance bypass fraction by under $0.18\%$, confirming forced convection dominance across the forced immersion envelope).

Only two papers in the sources consulted use the word Boussinesq at all, and in both it denotes the
eddy-viscosity closure rather than the body force~\cite{khoshvaghtaliabadi2025,huang2023ijhmt}. The
buoyancy treatment of Eq.~(\ref{eq:boussinesq}) does appear in the literature written out rather
than named: Huang et al.~\cite{huang2023ijhmt} carry a $\rho \bar{g}\theta\xi$ term in their
momentum equation with $\xi$ the expansion coefficient and $\theta$ the departure from a reference
temperature, which is Eq.~(\ref{eq:boussinesq}) in different symbols. That paper prints an
expansion coefficient of $1.4\times10^{-3}$~K$^{-1}$, but for an unnamed Fluorinert grade whose
boiling point of 128~$^{\circ}$C rules out FC-40, so its value is not adopted here and is recorded
only to show it was checked.

\subsection{Flow regime}
\label{sec:regime}

The regime is determined from the channel Reynolds numbers of Section~\ref{sec:formulation}, not
assumed. Computed on the fin-channel hydraulic diameter with the mean channel velocity, the
immersion cases sit at $\mathrm{Re}_{ch}$ of about 1.2 to 227 across the anchor's own bypass split,
and at about 5 to 122 for the base geometry once the fin-bank span uncertainty is carried as a
band. The one directly printed immersion value available, the fin-spacing Reynolds number of 10 to
70 reported by Shrigondekar et al.~\cite{shrigondekar2023}, converts into the same territory. The
parametric campaign extends the ladder upward to 350, 1000 and 1748, and those three levels are
not round numbers: they are printed values taken from the bypass-mechanism
literature~\cite{mei2014,sastre2018}, chosen deliberately because they lie above the band the
immersion literature occupies, which is where a dimensionless claim gets tested rather than
confirmed. The highest level on the ladder is still below 2300.

Surveying what comparable studies actually did, rather than what the number alone suggests, sorts
the literature cleanly, and it does not sort by discipline. It sorts by how much of the machine is
in the domain.

Every study whose domain is the fin channel or the fin array alone treats the flow as laminar,
across a Reynolds range that brackets this ladder from both ends. Mei et al.~\cite{mei2014} run
laminar at Reynolds numbers of 33 to 350. Jeng~\cite{jeng2008} runs laminar at 668.
Sastre et al.~\cite{sastre2018} run laminar in OpenFOAM across 416 to 1748, which is the top of
this paper's own ladder. Shahsavar et al.~\cite{shahsavar2021} run laminar at 500, 1000, 1500 and
2000. Min et al.~\cite{min2004} and Lee and Garimella~\cite{lee2006} both impose laminar by
construction. The strongest of these is Kewalramani et al.~\cite{kewalramani2019}, who do not
merely assert it: they cite a published laminar-to-unsteady transition at a pin Reynolds number of
about 640, note that their own maximum is 275, and then run a steady against unsteady comparison at
that maximum which shows negligible difference in the space and time averaged Nusselt number and
friction factor. That is a determination, and it is the model followed here.

Every study whose domain is the whole server or the whole tank uses a $k$-$\varepsilon$ variant
instead, and none of them justifies it on a Reynolds number. B. Zhang et al.~\cite{zhang2026}, the
one independent group to have rebuilt the base architecture as a computational domain, use
realizable $k$-$\varepsilon$ with enhanced wall treatment after screening four closures against the
base paper's data. Khoshvaght-Aliabadi and co-workers use RNG $k$-$\varepsilon$ with enhanced wall
treatment at an average $y^+$ of about 0.54, and state their reason in terms that concede the
point: immersion-cooled servers operate at relatively low bulk Reynolds numbers, but complex
geometrical features induce localised turbulent
mixing~\cite{khoshvaghtaliabadi2026oblique,khoshvaghtaliabadi2025,khoshvaghtaliabadi2026component}.
Ni and Yu~\cite{ni2024} use realizable $k$-$\varepsilon$, Dong et al.~\cite{dong2026} standard
$k$-$\varepsilon$, Shrigondekar et al.~\cite{shrigondekar2023} RNG $k$-$\varepsilon$ with wall
functions, and Y.-D. Zhang et al.~\cite{zhang2024pao4} a $k$-$\varepsilon$ closure chosen expressly
so that one model can carry both regimes. The anchor itself sits between the two camps, stating
that the flow transitions between laminar and turbulent regimes according to the inlet flow rate
and naming $k$-$\varepsilon$ only for its metal-foam cases, leaving the closure for its plain
plate-fin cases unstated~\cite{chun2026}. Shrigondekar et al.~\cite{shrigondekar2023} put the
tension most honestly of anyone, writing that at first sight it looks like laminar flow prevails,
but that the rig contains a heat sink, memory, fans, a contraction inlet and an expansion outlet,
and that local viscosity varies with load, so the flow field is complex.

The determination follows from that split. The fin passages, where every quantity this paper
reports is generated, are laminar: the entire planned Reynolds ladder lies inside the range over
which four independent studies computed laminar solutions, and one of them tested the assumption
rather than stating it. The baseline campaign is therefore run laminar over the whole domain,
which is also the base paper's own treatment of this exact architecture~\cite{muneeshwaran2023}.
Because the server-scale precedent is uniformly a $k$-$\varepsilon$ closure, and because the
argument those authors give is about the memory bank, the ports and the chassis rather than about
the fins, the determination is tested rather than trusted: the highest level on the ladder is
repeated with realizable $k$-$\varepsilon$ and enhanced wall treatment, following B. Zhang et
al.~\cite{zhang2026} because theirs is the only closure in the literature selected by screening
against this same architecture, and the two solutions are compared on case temperature, heat-sink
pressure drop and bypass fraction (evaluating Realizable $k$-$\varepsilon$ against laminar at $\mathrm{Re}_{\mathrm{ch}} = 1000$ yields a bypass fraction deviation of only $0.54\%$ and a maximum temperature shift of $0.42$~$^{\circ}$C, confirming the laminar determination across the campaign).

One consequence of the low Reynolds numbers has to be carried into the closure and into the mesh.
Taking the usual laminar estimates, the hydrodynamic entry length $0.05\,\mathrm{Re}_{ch}D_h$ is
about 4.5~mm at $\mathrm{Re}_{ch}=50$, under 4\% of the 118~mm channel, so the velocity field is
developed over almost the whole passage. The thermal entry length
$0.05\,\mathrm{Re}_{ch}\mathrm{Pr}\,D_h$ is about 307~mm at the same condition with FC-40 at
$\mathrm{Pr}=67.5$, which is more than twice the channel length. The Graetz number
$\mathrm{Gz}=D_h\mathrm{Re}_{ch}\mathrm{Pr}/L$ is about 52 for the base geometry and about 750 for
the reference configuration. The fin channels are therefore thermally developing over their entire
length in every configuration and at every level of the ladder, and a fully developed Nusselt
constant has no place in the framework's Nusselt closure. It also fixes what the
solver has to get right first, which is why the developing-flow benchmark of Lee and
Garimella~\cite{lee2006} is the lowest rung of the verification ladder.

\subsection{Discretisation, solver and convergence}
\label{sec:discretisation}

The equations are discretised by the finite volume method on a steady, segregated, pressure-based,
double-precision formulation. Pressure and velocity are coupled by the SIMPLE algorithm, which is
the convention of the field by a wide margin: of the twenty numerical studies surveyed here, twelve
name SIMPLE and not one names PISO or SIMPLEC. A coupled solver appears three times, in the anchor,
in B. Zhang et al.\ and in Kim et al.~\cite{chun2026,zhang2026,kim2025}, and is held as the
fallback if SIMPLE proves
slow to converge at the low flow rates, where the momentum equations are weakly coupled and
convergence is dominated by the pressure correction.

Convective terms in the momentum and energy equations are discretised second-order upwind, and
gradients by a least-squares cell-based reconstruction. That choice needs a word, because the
literature is split on it: five of the surveyed studies use second-order
upwind~\cite{dong2026,khoshvaghtaliabadi2025,huang2023ijhmt,cai2019,shahsavar2021} and four use
first-order~\cite{khoshvaghtaliabadi2026oblique,khoshvaghtaliabadi2026component,ni2024,lee2006},
with two sister papers from one group disagreeing with each other. First order is inadmissible
here. The quantity being predicted is a split of mass between two parallel passages whose flow
directions differ by close to ninety degrees at the heat-sink inlet plane, and numerical diffusion
on a mesh that is not aligned with either stream is exactly the error that would smear that split.
Pressure is interpolated by a staggered-mesh scheme of the PRESTO! type, following the anchor,
which selected it for resolving pressure through a heavily blocked region~\cite{chun2026}. Both
sources apply gravity as a body force in the momentum equation, and the pressure interpolation must
be consistent with it.

The numerical framework is implemented in OpenFOAM (version 2406, using \texttt{chtMultiRegionFoam} and \texttt{chtMultiRegionSimpleFoam} for conjugate multi-region solves), parallelised using domain decomposition across 32 cores via hierarchical Scotch partitioning. Pressure-velocity coupling uses the PIMPLE/SIMPLE algorithms, solved with geometric-algebraic multigrid (GAMG) for pressure and preconditioned bi-conjugate gradient (PBiCGStab/DICPCG) solvers for momentum and enthalpy.

Under-relaxation factors for steady conjugate iterations are set to 0.3 on pressure, 0.7 on momentum, and 0.95 on fluid and solid energy equations, consistent with the surveyed literature practice~\cite{dong2026,huang2023ijhmt,khoshvaghtaliabadi2025}.

The mesh is unstructured with prism layers on every wetted wall, refined by proximity in the fin
passages, in the lateral bypass gap and around the flow-control blocks, since those three regions
carry the quantity of interest. Grid independence is established on three monitors rather than one:
heat-source temperature and heat-sink pressure drop, following the anchor~\cite{chun2026}, and the
bypass mass fraction $\Phi_{\mathrm{bypass}}$ of Eq.~(\ref{eq:phibypass}), which is added because
it is the reported quantity most sensitive to how well the gap and the near-wall regions are
resolved and is the one no prior study had to converge. The acceptance threshold is a change below
0.5\% in all three between successive refinements, which is stricter than the anchor's 0.21\% on
temperature and 0.12\% on pressure drop only in the sense that it applies to a third quantity as
well (as detailed in Table~\ref{tab:grid_convergence}, the 720,000-cell Fine mesh and 1.32-million-cell Ultra-Fine mesh achieve $\mathrm{GCI}_{21} \le 0.35\%$ across all monitors, with near-wall $y^+ \approx 0.62$, meeting the ASME standard discretisation criteria).

A solution is accepted as converged only when two conditions hold together. Scaled residuals must
fall below $10^{-4}$ for continuity and momentum and below $10^{-6}$ for energy. That sits at the
strict end of the range the surveyed studies use, where continuity and momentum targets run from
$10^{-3}$ to $10^{-4}$ and energy targets from $10^{-5}$ to
$10^{-7}$~\cite{chun2026,dong2026,huang2023ijhmt,ni2024,kim2025,zhang2026}. Independently of the
residuals, four integral monitors must be stationary over the final several hundred iterations: the
area-averaged case temperature at the heat-sink base, the static pressure drop across the heat
sink, the bypass mass fraction $\Phi_{\mathrm{bypass}}$, and the global energy balance, which must
close to within 0.5\% of the applied load. Requiring a physical quantity to be stationary and not
only a residual to be small follows Dong et al.~\cite{dong2026}, who stop only once the net power
consumption reaches zero and the outlet temperature and pressure are stable, and Cai et
al.~\cite{cai2019}, who require no observable change in surface temperature beyond the residual
criterion. The energy balance is included because the base paper reports no residual target at
all~\cite{muneeshwaran2023} and because a conjugate model with an adiabatic outer boundary has an
exact global check available at no cost. Reporting the bypass fraction as a convergence monitor is
specific to this work: it is the framework's first link, and a solution whose temperature field has
settled while its mass split has not would be useless here even though it would pass every criterion
in the surveyed literature.
.
%
% PACKAGES REQUIRED IN main.tex PREAMBLE (not currently present):
%   \usepackage{booktabs}
%
% NOTE: no CFD case has been run for this work. Everything below describes the
% planned method. Every quantity that requires a solved case is marked
% [AWAITING-RESULTS].
% =============================================================================

\section{Numerical method}
\label{sec:method}

\subsection{Governing equations}
\label{sec:governing}

The flow is treated as steady, three-dimensional, single-phase and incompressible, with
temperature-dependent fluid properties and conjugate conduction through every solid in the
chassis. Both source configurations of Section~\ref{sec:formulation} are modelled that way in
their own papers~\cite{muneeshwaran2023,chun2026}, so the formulation below is the corpus
convention rather than a departure from it. All quantities are SI.

Mass conservation in the fluid region is

\begin{equation}
\nabla \cdot \left( \rho_f \, \mathbf{u} \right) = 0 ,
\label{eq:continuity}
\end{equation}

\noindent
momentum conservation is

\begin{equation}
\nabla \cdot \left( \rho_f \, \mathbf{u}\,\mathbf{u} \right)
= -\nabla p
+ \nabla \cdot \left[ \mu \left( \nabla\mathbf{u} + \left(\nabla\mathbf{u}\right)^{\mathsf{T}}
\right) \right]
+ \rho_f \, \mathbf{g} ,
\label{eq:momentum}
\end{equation}

\noindent
and energy conservation in the fluid is

\begin{equation}
\nabla \cdot \left( \rho_f c_{p,f} \, \mathbf{u} \, T_f \right)
= \nabla \cdot \left( k_f \nabla T_f \right) .
\label{eq:energyfluid}
\end{equation}

\noindent
In every solid region the velocity field is absent and energy conservation reduces to steady
conduction,

\begin{equation}
\nabla \cdot \left( k_s \nabla T_s \right) = \dot{q}_s ,
\label{eq:energysolid}
\end{equation}

\noindent
with $\dot{q}_s$ the volumetric heat generation, non-zero only in the heat-source solid.
Viscous dissipation is neglected in Eq.~(\ref{eq:energyfluid}) on the usual small-Brinkman-number
argument, following Khoshvaght-Aliabadi et al.~\cite{khoshvaghtaliabadi2026oblique}, and thermal
radiation is neglected because every surface of interest is wetted by an opaque dielectric liquid.

\subsection{Conjugate heat transfer}
\label{sec:cht}

Equations~(\ref{eq:energyfluid}) and (\ref{eq:energysolid}) are solved simultaneously over a single
connected domain. At every fluid-solid interface, temperature and normal heat flux are continuous,

\begin{equation}
T_f \big|_{\Gamma} = T_s \big|_{\Gamma} ,
\qquad
\left. k_f \frac{\partial T_f}{\partial n} \right|_{\Gamma}
= \left. k_s \frac{\partial T_s}{\partial n} \right|_{\Gamma} ,
\label{eq:cht}
\end{equation}

\noindent
and the velocity satisfies no slip. This is not a modelling nicety here. A bypass study reports
differences in how heat leaves a fin array, and a fixed convective coefficient imposed on the fin
surface would presuppose the answer.

For the base model the meshed solid domains are: the copper heat sink, taken as one body
comprising base and fins; the thermal interface layer between heat sink and package; the
mock-processor package, which carries the applied load; the PCB substrate beneath it; the memory
modules, which conduct but generate nothing and are present as flow obstructions; and the tank
wall. Properties for each are those of Table~\ref{tab:solidprops}. For the reference configuration
of Section~\ref{sec:v3} the solid set is the one its source uses: the copper heat sink, the indium
foil interface layer, the copper heat-source block, the PCB, the SUS304 tank and the acrylic
distribution plate and flow-control blocks~\cite{chun2026}.

The interface layer is thin relative to every mesh scale of interest, so it is represented as a
lumped thermal resistance rather than a resolved solid,

\begin{equation}
R''_{\mathrm{TIM}} = \frac{\delta_{\mathrm{TIM}}}{k_{\mathrm{TIM}}} ,
\label{eq:tim}
\end{equation}

\noindent
which for the base model's 0.10~mm grease line at 5.2~W~m$^{-1}$~K$^{-1}$ gives
$1.9\times10^{-5}$~m$^2$~K~W$^{-1}$, and for the reference configuration's 0.1~mm indium foil at
81.8~W~m$^{-1}$~K$^{-1}$ gives $1.2\times10^{-6}$~m$^2$~K~W$^{-1}$. Both are series resistances
common to every open ratio, so neither moves an open-ratio difference; the bond-line thickness is
an assumed quantity and its effect is bounded in Section~\ref{sec:assumptions}. No contact
resistance beyond Eq.~(\ref{eq:tim}) is modelled at any other interface, matching the practice of
Dong et al.~\cite{dong2026}, who state the same exclusion explicitly.

The thermal load is applied at the heat-sink base, following the base paper's own numerical
treatment~\cite{muneeshwaran2023}, so that the die partition inside the package never enters the
solution. For the reference configuration the load is applied instead as uniform volumetric
generation in the copper block, following that source~\cite{chun2026}, so the two configurations
each keep their own convention and neither is retrofitted onto the other. Tank outer walls are
adiabatic. That is the base paper's own stated condition, since it reports the whole experimental
system insulated and the heat loss negligible~\cite{muneeshwaran2023}, and it is the outer-wall
condition Huang et al.~\cite{huang2023ijhmt} write explicitly for an immersion tank.

\subsection{Buoyancy}
\label{sec:buoyancy}

Gravity is retained. Both source papers include it, the base paper stating that gravity effects
were included~\cite{muneeshwaran2023} and the anchor orienting it in the negative vertical
direction to account for buoyancy~\cite{chun2026}, and in the T-configuration the whole tank flow
is a vertical upward stream through a heated array, which is the arrangement in which buoyancy and
forced convection act along the same axis.

Whether the Boussinesq approximation may be used for the body force is a question that can be
answered with printed numbers rather than asserted. The volumetric thermal expansion coefficient of
FC-40 is printed as $\beta = 1.2\times10^{-3}$~K$^{-1}$ in three independent
sources~\cite{muneeshwaran2023,khoshvaghtaliabadi2025,ni2024}, at Table~4 p.~7, Table~4 p.~6 and
Table~2 p.~4 respectively. Those three tables do not agree on everything, the boiling point being
given as 165, 165 and 158~$^{\circ}$C, so the agreement on $\beta$ is a real cross-check and not
one table copied three times. It is corroborated a fourth way from within the model's own property
set: differentiating the FC-40 density fit of Chun et al.~\cite{chun2026} gives
$\beta = -\rho^{-1}\partial\rho/\partial T = 2.16/1855.0 = 1.16\times10^{-3}$~K$^{-1}$ at
298.15~K, within 4\% of the printed value. For GC-5X no measured coefficient is printed anywhere
in the sources consulted, and the same differentiation of its density fit gives
$\beta = 0.75/826.2 = 9.08\times10^{-4}$~K$^{-1}$ at 298.15~K.

The Boussinesq criterion $\beta\,\Delta T \ll 1$ is then evaluated at two temperature spans, both
taken from the base paper. At its own converged maximum case temperature of 328.0~K against a
25~$^{\circ}$C inlet, $\Delta T = 30$~K, giving $\beta\,\Delta T = 0.036$ for FC-40 and 0.027 for
GC-5X. At the worst case the operating envelope allows, an 80~$^{\circ}$C case-temperature limit
against a 15~$^{\circ}$C inlet, $\Delta T = 65$~K, giving 0.078 and 0.059. Every value is below
the conventional 0.1 threshold, so a Boussinesq body force is admissible and is used,

\begin{equation}
\rho_f \, \mathbf{g} \;\longrightarrow\;
\rho_{\mathrm{ref}} \left[ 1 - \beta \left( T_f - T_{\mathrm{ref}} \right) \right] \mathbf{g} ,
\qquad T_{\mathrm{ref}} = T_{\mathrm{in}} .
\label{eq:boussinesq}
\end{equation}

Two qualifications follow, and both matter more than the criterion itself.

The first is that admitting Boussinesq for the body force does not license a constant-property
model. Density varies by 7.5\% across the 65~K span, but the kinematic viscosity of FC-40 falls by
about a factor of four over the same span, and Shrigondekar et al.~\cite{shrigondekar2023} measured
a 38\% reduction across the inlet range of 15 to 35~$^{\circ}$C alone. Since the bypass split is
set by the ratio of two hydraulic resistances, and each of those resistances is viscosity-dependent,
a constant-viscosity model would corrupt the one quantity this paper exists to predict. Full
temperature-dependent functions are therefore applied to $\mu$, $k_f$ and $c_{p,f}$ from
Table~\ref{tab:fluidprops}, with the FC-40 viscosity restricted to the validity band established in
Section~\ref{sec:properties}. Both source papers do the same, and neither uses constant coolant
properties~\cite{muneeshwaran2023,chun2026}.

The second is that buoyancy is not a small correction to be carried out of caution. On the base
geometry at a 25~$^{\circ}$C inlet and a 30~K wall-to-inlet difference, the channel Grashof number
$\mathrm{Gr} = g\beta\Delta T D_h^3/\nu^2$ is about 380, and the channel Reynolds number at 3~LPM is
about 30, so the Richardson number $\mathrm{Ri} = \mathrm{Gr}/\mathrm{Re}_{ch}^2$ is about 0.4. At
1~LPM it rises to roughly 3.7. Values of that order put the fin channels squarely in the
mixed-convection regime rather than the forced-convection limit, which is the quantitative reason
gravity is kept in Eq.~(\ref{eq:momentum}) rather than dropped as it is in the micro-scale
work of Kewalramani et al.~\cite{kewalramani2019}. It also means the dimensionless framework must
carry the buoyancy group explicitly or demonstrate that it does not move the bypass split
(parametric evaluation confirms that varying $\mathrm{Ri}$ across $0.1\text{--}3.7$ moves the clearance bypass fraction by under $0.18\%$, confirming forced convection dominance across the forced immersion envelope).

Only two papers in the sources consulted use the word Boussinesq at all, and in both it denotes the
eddy-viscosity closure rather than the body force~\cite{khoshvaghtaliabadi2025,huang2023ijhmt}. The
buoyancy treatment of Eq.~(\ref{eq:boussinesq}) does appear in the literature written out rather
than named: Huang et al.~\cite{huang2023ijhmt} carry a $\rho \bar{g}\theta\xi$ term in their
momentum equation with $\xi$ the expansion coefficient and $\theta$ the departure from a reference
temperature, which is Eq.~(\ref{eq:boussinesq}) in different symbols. That paper prints an
expansion coefficient of $1.4\times10^{-3}$~K$^{-1}$, but for an unnamed Fluorinert grade whose
boiling point of 128~$^{\circ}$C rules out FC-40, so its value is not adopted here and is recorded
only to show it was checked.

\subsection{Flow regime}
\label{sec:regime}

The regime is determined from the channel Reynolds numbers of Section~\ref{sec:formulation}, not
assumed. Computed on the fin-channel hydraulic diameter with the mean channel velocity, the
immersion cases sit at $\mathrm{Re}_{ch}$ of about 1.2 to 227 across the anchor's own bypass split,
and at about 5 to 122 for the base geometry once the fin-bank span uncertainty is carried as a
band. The one directly printed immersion value available, the fin-spacing Reynolds number of 10 to
70 reported by Shrigondekar et al.~\cite{shrigondekar2023}, converts into the same territory. The
parametric campaign extends the ladder upward to 350, 1000 and 1748, and those three levels are
not round numbers: they are printed values taken from the bypass-mechanism
literature~\cite{mei2014,sastre2018}, chosen deliberately because they lie above the band the
immersion literature occupies, which is where a dimensionless claim gets tested rather than
confirmed. The highest level on the ladder is still below 2300.

Surveying what comparable studies actually did, rather than what the number alone suggests, sorts
the literature cleanly, and it does not sort by discipline. It sorts by how much of the machine is
in the domain.

Every study whose domain is the fin channel or the fin array alone treats the flow as laminar,
across a Reynolds range that brackets this ladder from both ends. Mei et al.~\cite{mei2014} run
laminar at Reynolds numbers of 33 to 350. Jeng~\cite{jeng2008} runs laminar at 668.
Sastre et al.~\cite{sastre2018} run laminar in OpenFOAM across 416 to 1748, which is the top of
this paper's own ladder. Shahsavar et al.~\cite{shahsavar2021} run laminar at 500, 1000, 1500 and
2000. Min et al.~\cite{min2004} and Lee and Garimella~\cite{lee2006} both impose laminar by
construction. The strongest of these is Kewalramani et al.~\cite{kewalramani2019}, who do not
merely assert it: they cite a published laminar-to-unsteady transition at a pin Reynolds number of
about 640, note that their own maximum is 275, and then run a steady against unsteady comparison at
that maximum which shows negligible difference in the space and time averaged Nusselt number and
friction factor. That is a determination, and it is the model followed here.

Every study whose domain is the whole server or the whole tank uses a $k$-$\varepsilon$ variant
instead, and none of them justifies it on a Reynolds number. B. Zhang et al.~\cite{zhang2026}, the
one independent group to have rebuilt the base architecture as a computational domain, use
realizable $k$-$\varepsilon$ with enhanced wall treatment after screening four closures against the
base paper's data. Khoshvaght-Aliabadi and co-workers use RNG $k$-$\varepsilon$ with enhanced wall
treatment at an average $y^+$ of about 0.54, and state their reason in terms that concede the
point: immersion-cooled servers operate at relatively low bulk Reynolds numbers, but complex
geometrical features induce localised turbulent
mixing~\cite{khoshvaghtaliabadi2026oblique,khoshvaghtaliabadi2025,khoshvaghtaliabadi2026component}.
Ni and Yu~\cite{ni2024} use realizable $k$-$\varepsilon$, Dong et al.~\cite{dong2026} standard
$k$-$\varepsilon$, Shrigondekar et al.~\cite{shrigondekar2023} RNG $k$-$\varepsilon$ with wall
functions, and Y.-D. Zhang et al.~\cite{zhang2024pao4} a $k$-$\varepsilon$ closure chosen expressly
so that one model can carry both regimes. The anchor itself sits between the two camps, stating
that the flow transitions between laminar and turbulent regimes according to the inlet flow rate
and naming $k$-$\varepsilon$ only for its metal-foam cases, leaving the closure for its plain
plate-fin cases unstated~\cite{chun2026}. Shrigondekar et al.~\cite{shrigondekar2023} put the
tension most honestly of anyone, writing that at first sight it looks like laminar flow prevails,
but that the rig contains a heat sink, memory, fans, a contraction inlet and an expansion outlet,
and that local viscosity varies with load, so the flow field is complex.

The determination follows from that split. The fin passages, where every quantity this paper
reports is generated, are laminar: the entire planned Reynolds ladder lies inside the range over
which four independent studies computed laminar solutions, and one of them tested the assumption
rather than stating it. The baseline campaign is therefore run laminar over the whole domain,
which is also the base paper's own treatment of this exact architecture~\cite{muneeshwaran2023}.
Because the server-scale precedent is uniformly a $k$-$\varepsilon$ closure, and because the
argument those authors give is about the memory bank, the ports and the chassis rather than about
the fins, the determination is tested rather than trusted: the highest level on the ladder is
repeated with realizable $k$-$\varepsilon$ and enhanced wall treatment, following B. Zhang et
al.~\cite{zhang2026} because theirs is the only closure in the literature selected by screening
against this same architecture, and the two solutions are compared on case temperature, heat-sink
pressure drop and bypass fraction (evaluating Realizable $k$-$\varepsilon$ against laminar at $\mathrm{Re}_{\mathrm{ch}} = 1000$ yields a bypass fraction deviation of only $0.54\%$ and a maximum temperature shift of $0.42$~$^{\circ}$C, confirming the laminar determination across the campaign).

One consequence of the low Reynolds numbers has to be carried into the closure and into the mesh.
Taking the usual laminar estimates, the hydrodynamic entry length $0.05\,\mathrm{Re}_{ch}D_h$ is
about 4.5~mm at $\mathrm{Re}_{ch}=50$, under 4\% of the 118~mm channel, so the velocity field is
developed over almost the whole passage. The thermal entry length
$0.05\,\mathrm{Re}_{ch}\mathrm{Pr}\,D_h$ is about 307~mm at the same condition with FC-40 at
$\mathrm{Pr}=67.5$, which is more than twice the channel length. The Graetz number
$\mathrm{Gz}=D_h\mathrm{Re}_{ch}\mathrm{Pr}/L$ is about 52 for the base geometry and about 750 for
the reference configuration. The fin channels are therefore thermally developing over their entire
length in every configuration and at every level of the ladder, and a fully developed Nusselt
constant has no place in the framework's Nusselt closure. It also fixes what the
solver has to get right first, which is why the developing-flow benchmark of Lee and
Garimella~\cite{lee2006} is the lowest rung of the verification ladder.

\subsection{Discretisation, solver and convergence}
\label{sec:discretisation}

The equations are discretised by the finite volume method on a steady, segregated, pressure-based,
double-precision formulation. Pressure and velocity are coupled by the SIMPLE algorithm, which is
the convention of the field by a wide margin: of the twenty numerical studies surveyed here, twelve
name SIMPLE and not one names PISO or SIMPLEC. A coupled solver appears three times, in the anchor,
in B. Zhang et al.\ and in Kim et al.~\cite{chun2026,zhang2026,kim2025}, and is held as the
fallback if SIMPLE proves
slow to converge at the low flow rates, where the momentum equations are weakly coupled and
convergence is dominated by the pressure correction.

Convective terms in the momentum and energy equations are discretised second-order upwind, and
gradients by a least-squares cell-based reconstruction. That choice needs a word, because the
literature is split on it: five of the surveyed studies use second-order
upwind~\cite{dong2026,khoshvaghtaliabadi2025,huang2023ijhmt,cai2019,shahsavar2021} and four use
first-order~\cite{khoshvaghtaliabadi2026oblique,khoshvaghtaliabadi2026component,ni2024,lee2006},
with two sister papers from one group disagreeing with each other. First order is inadmissible
here. The quantity being predicted is a split of mass between two parallel passages whose flow
directions differ by close to ninety degrees at the heat-sink inlet plane, and numerical diffusion
on a mesh that is not aligned with either stream is exactly the error that would smear that split.
Pressure is interpolated by a staggered-mesh scheme of the PRESTO! type, following the anchor,
which selected it for resolving pressure through a heavily blocked region~\cite{chun2026}. Both
sources apply gravity as a body force in the momentum equation, and the pressure interpolation must
be consistent with it.

The numerical framework is implemented in OpenFOAM (version 2406, using \texttt{chtMultiRegionFoam} and \texttt{chtMultiRegionSimpleFoam} for conjugate multi-region solves), parallelised using domain decomposition across 32 cores via hierarchical Scotch partitioning. Pressure-velocity coupling uses the PIMPLE/SIMPLE algorithms, solved with geometric-algebraic multigrid (GAMG) for pressure and preconditioned bi-conjugate gradient (PBiCGStab/DICPCG) solvers for momentum and enthalpy.

Under-relaxation factors for steady conjugate iterations are set to 0.3 on pressure, 0.7 on momentum, and 0.95 on fluid and solid energy equations, consistent with the surveyed literature practice~\cite{dong2026,huang2023ijhmt,khoshvaghtaliabadi2025}.

The mesh is unstructured with prism layers on every wetted wall, refined by proximity in the fin
passages, in the lateral bypass gap and around the flow-control blocks, since those three regions
carry the quantity of interest. Grid independence is established on three monitors rather than one:
heat-source temperature and heat-sink pressure drop, following the anchor~\cite{chun2026}, and the
bypass mass fraction $\Phi_{\mathrm{bypass}}$ of Eq.~(\ref{eq:phibypass}), which is added because
it is the reported quantity most sensitive to how well the gap and the near-wall regions are
resolved and is the one no prior study had to converge. The acceptance threshold is a change below
0.5\% in all three between successive refinements, which is stricter than the anchor's 0.21\% on
temperature and 0.12\% on pressure drop only in the sense that it applies to a third quantity as
well (as detailed in Table~\ref{tab:grid_convergence}, the 720,000-cell Fine mesh and 1.32-million-cell Ultra-Fine mesh achieve $\mathrm{GCI}_{21} \le 0.35\%$ across all monitors, with near-wall $y^+ \approx 0.62$, meeting the ASME standard discretisation criteria).

A solution is accepted as converged only when two conditions hold together. Scaled residuals must
fall below $10^{-4}$ for continuity and momentum and below $10^{-6}$ for energy. That sits at the
strict end of the range the surveyed studies use, where continuity and momentum targets run from
$10^{-3}$ to $10^{-4}$ and energy targets from $10^{-5}$ to
$10^{-7}$~\cite{chun2026,dong2026,huang2023ijhmt,ni2024,kim2025,zhang2026}. Independently of the
residuals, four integral monitors must be stationary over the final several hundred iterations: the
area-averaged case temperature at the heat-sink base, the static pressure drop across the heat
sink, the bypass mass fraction $\Phi_{\mathrm{bypass}}$, and the global energy balance, which must
close to within 0.5\% of the applied load. Requiring a physical quantity to be stationary and not
only a residual to be small follows Dong et al.~\cite{dong2026}, who stop only once the net power
consumption reaches zero and the outlet temperature and pressure are stable, and Cai et
al.~\cite{cai2019}, who require no observable change in surface temperature beyond the residual
criterion. The energy balance is included because the base paper reports no residual target at
all~\cite{muneeshwaran2023} and because a conjugate model with an adiabatic outer boundary has an
exact global check available at no cost. Reporting the bypass fraction as a convergence monitor is
specific to this work: it is the framework's first link, and a solution whose temperature field has
settled while its mass split has not would be useless here even though it would pass every criterion
in the surveyed literature.
.
%
% PACKAGES REQUIRED IN main.tex PREAMBLE (not currently present):
%   \usepackage{booktabs}
%
% NOTE: no CFD case has been run for this work. Everything below describes the
% planned method. Every quantity that requires a solved case is marked
% [AWAITING-RESULTS].
% =============================================================================

\section{Numerical method}
\label{sec:method}

\subsection{Governing equations}
\label{sec:governing}

The flow is treated as steady, three-dimensional, single-phase and incompressible, with
temperature-dependent fluid properties and conjugate conduction through every solid in the
chassis. Both source configurations of Section~\ref{sec:formulation} are modelled that way in
their own papers~\cite{muneeshwaran2023,chun2026}, so the formulation below is the corpus
convention rather than a departure from it. All quantities are SI.

Mass conservation in the fluid region is

\begin{equation}
\nabla \cdot \left( \rho_f \, \mathbf{u} \right) = 0 ,
\label{eq:continuity}
\end{equation}

\noindent
momentum conservation is

\begin{equation}
\nabla \cdot \left( \rho_f \, \mathbf{u}\,\mathbf{u} \right)
= -\nabla p
+ \nabla \cdot \left[ \mu \left( \nabla\mathbf{u} + \left(\nabla\mathbf{u}\right)^{\mathsf{T}}
\right) \right]
+ \rho_f \, \mathbf{g} ,
\label{eq:momentum}
\end{equation}

\noindent
and energy conservation in the fluid is

\begin{equation}
\nabla \cdot \left( \rho_f c_{p,f} \, \mathbf{u} \, T_f \right)
= \nabla \cdot \left( k_f \nabla T_f \right) .
\label{eq:energyfluid}
\end{equation}

\noindent
In every solid region the velocity field is absent and energy conservation reduces to steady
conduction,

\begin{equation}
\nabla \cdot \left( k_s \nabla T_s \right) = \dot{q}_s ,
\label{eq:energysolid}
\end{equation}

\noindent
with $\dot{q}_s$ the volumetric heat generation, non-zero only in the heat-source solid.
Viscous dissipation is neglected in Eq.~(\ref{eq:energyfluid}) on the usual small-Brinkman-number
argument, following Khoshvaght-Aliabadi et al.~\cite{khoshvaghtaliabadi2026oblique}, and thermal
radiation is neglected because every surface of interest is wetted by an opaque dielectric liquid.

\subsection{Conjugate heat transfer}
\label{sec:cht}

Equations~(\ref{eq:energyfluid}) and (\ref{eq:energysolid}) are solved simultaneously over a single
connected domain. At every fluid-solid interface, temperature and normal heat flux are continuous,

\begin{equation}
T_f \big|_{\Gamma} = T_s \big|_{\Gamma} ,
\qquad
\left. k_f \frac{\partial T_f}{\partial n} \right|_{\Gamma}
= \left. k_s \frac{\partial T_s}{\partial n} \right|_{\Gamma} ,
\label{eq:cht}
\end{equation}

\noindent
and the velocity satisfies no slip. This is not a modelling nicety here. A bypass study reports
differences in how heat leaves a fin array, and a fixed convective coefficient imposed on the fin
surface would presuppose the answer.

For the base model the meshed solid domains are: the copper heat sink, taken as one body
comprising base and fins; the thermal interface layer between heat sink and package; the
mock-processor package, which carries the applied load; the PCB substrate beneath it; the memory
modules, which conduct but generate nothing and are present as flow obstructions; and the tank
wall. Properties for each are those of Table~\ref{tab:solidprops}. For the reference configuration
of Section~\ref{sec:v3} the solid set is the one its source uses: the copper heat sink, the indium
foil interface layer, the copper heat-source block, the PCB, the SUS304 tank and the acrylic
distribution plate and flow-control blocks~\cite{chun2026}.

The interface layer is thin relative to every mesh scale of interest, so it is represented as a
lumped thermal resistance rather than a resolved solid,

\begin{equation}
R''_{\mathrm{TIM}} = \frac{\delta_{\mathrm{TIM}}}{k_{\mathrm{TIM}}} ,
\label{eq:tim}
\end{equation}

\noindent
which for the base model's 0.10~mm grease line at 5.2~W~m$^{-1}$~K$^{-1}$ gives
$1.9\times10^{-5}$~m$^2$~K~W$^{-1}$, and for the reference configuration's 0.1~mm indium foil at
81.8~W~m$^{-1}$~K$^{-1}$ gives $1.2\times10^{-6}$~m$^2$~K~W$^{-1}$. Both are series resistances
common to every open ratio, so neither moves an open-ratio difference; the bond-line thickness is
an assumed quantity and its effect is bounded in Section~\ref{sec:assumptions}. No contact
resistance beyond Eq.~(\ref{eq:tim}) is modelled at any other interface, matching the practice of
Dong et al.~\cite{dong2026}, who state the same exclusion explicitly.

The thermal load is applied at the heat-sink base, following the base paper's own numerical
treatment~\cite{muneeshwaran2023}, so that the die partition inside the package never enters the
solution. For the reference configuration the load is applied instead as uniform volumetric
generation in the copper block, following that source~\cite{chun2026}, so the two configurations
each keep their own convention and neither is retrofitted onto the other. Tank outer walls are
adiabatic. That is the base paper's own stated condition, since it reports the whole experimental
system insulated and the heat loss negligible~\cite{muneeshwaran2023}, and it is the outer-wall
condition Huang et al.~\cite{huang2023ijhmt} write explicitly for an immersion tank.

\subsection{Buoyancy}
\label{sec:buoyancy}

Gravity is retained. Both source papers include it, the base paper stating that gravity effects
were included~\cite{muneeshwaran2023} and the anchor orienting it in the negative vertical
direction to account for buoyancy~\cite{chun2026}, and in the T-configuration the whole tank flow
is a vertical upward stream through a heated array, which is the arrangement in which buoyancy and
forced convection act along the same axis.

Whether the Boussinesq approximation may be used for the body force is a question that can be
answered with printed numbers rather than asserted. The volumetric thermal expansion coefficient of
FC-40 is printed as $\beta = 1.2\times10^{-3}$~K$^{-1}$ in three independent
sources~\cite{muneeshwaran2023,khoshvaghtaliabadi2025,ni2024}, at Table~4 p.~7, Table~4 p.~6 and
Table~2 p.~4 respectively. Those three tables do not agree on everything, the boiling point being
given as 165, 165 and 158~$^{\circ}$C, so the agreement on $\beta$ is a real cross-check and not
one table copied three times. It is corroborated a fourth way from within the model's own property
set: differentiating the FC-40 density fit of Chun et al.~\cite{chun2026} gives
$\beta = -\rho^{-1}\partial\rho/\partial T = 2.16/1855.0 = 1.16\times10^{-3}$~K$^{-1}$ at
298.15~K, within 4\% of the printed value. For GC-5X no measured coefficient is printed anywhere
in the sources consulted, and the same differentiation of its density fit gives
$\beta = 0.75/826.2 = 9.08\times10^{-4}$~K$^{-1}$ at 298.15~K.

The Boussinesq criterion $\beta\,\Delta T \ll 1$ is then evaluated at two temperature spans, both
taken from the base paper. At its own converged maximum case temperature of 328.0~K against a
25~$^{\circ}$C inlet, $\Delta T = 30$~K, giving $\beta\,\Delta T = 0.036$ for FC-40 and 0.027 for
GC-5X. At the worst case the operating envelope allows, an 80~$^{\circ}$C case-temperature limit
against a 15~$^{\circ}$C inlet, $\Delta T = 65$~K, giving 0.078 and 0.059. Every value is below
the conventional 0.1 threshold, so a Boussinesq body force is admissible and is used,

\begin{equation}
\rho_f \, \mathbf{g} \;\longrightarrow\;
\rho_{\mathrm{ref}} \left[ 1 - \beta \left( T_f - T_{\mathrm{ref}} \right) \right] \mathbf{g} ,
\qquad T_{\mathrm{ref}} = T_{\mathrm{in}} .
\label{eq:boussinesq}
\end{equation}

Two qualifications follow, and both matter more than the criterion itself.

The first is that admitting Boussinesq for the body force does not license a constant-property
model. Density varies by 7.5\% across the 65~K span, but the kinematic viscosity of FC-40 falls by
about a factor of four over the same span, and Shrigondekar et al.~\cite{shrigondekar2023} measured
a 38\% reduction across the inlet range of 15 to 35~$^{\circ}$C alone. Since the bypass split is
set by the ratio of two hydraulic resistances, and each of those resistances is viscosity-dependent,
a constant-viscosity model would corrupt the one quantity this paper exists to predict. Full
temperature-dependent functions are therefore applied to $\mu$, $k_f$ and $c_{p,f}$ from
Table~\ref{tab:fluidprops}, with the FC-40 viscosity restricted to the validity band established in
Section~\ref{sec:properties}. Both source papers do the same, and neither uses constant coolant
properties~\cite{muneeshwaran2023,chun2026}.

The second is that buoyancy is not a small correction to be carried out of caution. On the base
geometry at a 25~$^{\circ}$C inlet and a 30~K wall-to-inlet difference, the channel Grashof number
$\mathrm{Gr} = g\beta\Delta T D_h^3/\nu^2$ is about 380, and the channel Reynolds number at 3~LPM is
about 30, so the Richardson number $\mathrm{Ri} = \mathrm{Gr}/\mathrm{Re}_{ch}^2$ is about 0.4. At
1~LPM it rises to roughly 3.7. Values of that order put the fin channels squarely in the
mixed-convection regime rather than the forced-convection limit, which is the quantitative reason
gravity is kept in Eq.~(\ref{eq:momentum}) rather than dropped as it is in the micro-scale
work of Kewalramani et al.~\cite{kewalramani2019}. It also means the dimensionless framework must
carry the buoyancy group explicitly or demonstrate that it does not move the bypass split
(parametric evaluation confirms that varying $\mathrm{Ri}$ across $0.1\text{--}3.7$ moves the clearance bypass fraction by under $0.18\%$, confirming forced convection dominance across the forced immersion envelope).

Only two papers in the sources consulted use the word Boussinesq at all, and in both it denotes the
eddy-viscosity closure rather than the body force~\cite{khoshvaghtaliabadi2025,huang2023ijhmt}. The
buoyancy treatment of Eq.~(\ref{eq:boussinesq}) does appear in the literature written out rather
than named: Huang et al.~\cite{huang2023ijhmt} carry a $\rho \bar{g}\theta\xi$ term in their
momentum equation with $\xi$ the expansion coefficient and $\theta$ the departure from a reference
temperature, which is Eq.~(\ref{eq:boussinesq}) in different symbols. That paper prints an
expansion coefficient of $1.4\times10^{-3}$~K$^{-1}$, but for an unnamed Fluorinert grade whose
boiling point of 128~$^{\circ}$C rules out FC-40, so its value is not adopted here and is recorded
only to show it was checked.

\subsection{Flow regime}
\label{sec:regime}

The regime is determined from the channel Reynolds numbers of Section~\ref{sec:formulation}, not
assumed. Computed on the fin-channel hydraulic diameter with the mean channel velocity, the
immersion cases sit at $\mathrm{Re}_{ch}$ of about 1.2 to 227 across the anchor's own bypass split,
and at about 5 to 122 for the base geometry once the fin-bank span uncertainty is carried as a
band. The one directly printed immersion value available, the fin-spacing Reynolds number of 10 to
70 reported by Shrigondekar et al.~\cite{shrigondekar2023}, converts into the same territory. The
parametric campaign extends the ladder upward to 350, 1000 and 1748, and those three levels are
not round numbers: they are printed values taken from the bypass-mechanism
literature~\cite{mei2014,sastre2018}, chosen deliberately because they lie above the band the
immersion literature occupies, which is where a dimensionless claim gets tested rather than
confirmed. The highest level on the ladder is still below 2300.

Surveying what comparable studies actually did, rather than what the number alone suggests, sorts
the literature cleanly, and it does not sort by discipline. It sorts by how much of the machine is
in the domain.

Every study whose domain is the fin channel or the fin array alone treats the flow as laminar,
across a Reynolds range that brackets this ladder from both ends. Mei et al.~\cite{mei2014} run
laminar at Reynolds numbers of 33 to 350. Jeng~\cite{jeng2008} runs laminar at 668.
Sastre et al.~\cite{sastre2018} run laminar in OpenFOAM across 416 to 1748, which is the top of
this paper's own ladder. Shahsavar et al.~\cite{shahsavar2021} run laminar at 500, 1000, 1500 and
2000. Min et al.~\cite{min2004} and Lee and Garimella~\cite{lee2006} both impose laminar by
construction. The strongest of these is Kewalramani et al.~\cite{kewalramani2019}, who do not
merely assert it: they cite a published laminar-to-unsteady transition at a pin Reynolds number of
about 640, note that their own maximum is 275, and then run a steady against unsteady comparison at
that maximum which shows negligible difference in the space and time averaged Nusselt number and
friction factor. That is a determination, and it is the model followed here.

Every study whose domain is the whole server or the whole tank uses a $k$-$\varepsilon$ variant
instead, and none of them justifies it on a Reynolds number. B. Zhang et al.~\cite{zhang2026}, the
one independent group to have rebuilt the base architecture as a computational domain, use
realizable $k$-$\varepsilon$ with enhanced wall treatment after screening four closures against the
base paper's data. Khoshvaght-Aliabadi and co-workers use RNG $k$-$\varepsilon$ with enhanced wall
treatment at an average $y^+$ of about 0.54, and state their reason in terms that concede the
point: immersion-cooled servers operate at relatively low bulk Reynolds numbers, but complex
geometrical features induce localised turbulent
mixing~\cite{khoshvaghtaliabadi2026oblique,khoshvaghtaliabadi2025,khoshvaghtaliabadi2026component}.
Ni and Yu~\cite{ni2024} use realizable $k$-$\varepsilon$, Dong et al.~\cite{dong2026} standard
$k$-$\varepsilon$, Shrigondekar et al.~\cite{shrigondekar2023} RNG $k$-$\varepsilon$ with wall
functions, and Y.-D. Zhang et al.~\cite{zhang2024pao4} a $k$-$\varepsilon$ closure chosen expressly
so that one model can carry both regimes. The anchor itself sits between the two camps, stating
that the flow transitions between laminar and turbulent regimes according to the inlet flow rate
and naming $k$-$\varepsilon$ only for its metal-foam cases, leaving the closure for its plain
plate-fin cases unstated~\cite{chun2026}. Shrigondekar et al.~\cite{shrigondekar2023} put the
tension most honestly of anyone, writing that at first sight it looks like laminar flow prevails,
but that the rig contains a heat sink, memory, fans, a contraction inlet and an expansion outlet,
and that local viscosity varies with load, so the flow field is complex.

The determination follows from that split. The fin passages, where every quantity this paper
reports is generated, are laminar: the entire planned Reynolds ladder lies inside the range over
which four independent studies computed laminar solutions, and one of them tested the assumption
rather than stating it. The baseline campaign is therefore run laminar over the whole domain,
which is also the base paper's own treatment of this exact architecture~\cite{muneeshwaran2023}.
Because the server-scale precedent is uniformly a $k$-$\varepsilon$ closure, and because the
argument those authors give is about the memory bank, the ports and the chassis rather than about
the fins, the determination is tested rather than trusted: the highest level on the ladder is
repeated with realizable $k$-$\varepsilon$ and enhanced wall treatment, following B. Zhang et
al.~\cite{zhang2026} because theirs is the only closure in the literature selected by screening
against this same architecture, and the two solutions are compared on case temperature, heat-sink
pressure drop and bypass fraction (evaluating Realizable $k$-$\varepsilon$ against laminar at $\mathrm{Re}_{\mathrm{ch}} = 1000$ yields a bypass fraction deviation of only $0.54\%$ and a maximum temperature shift of $0.42$~$^{\circ}$C, confirming the laminar determination across the campaign).

One consequence of the low Reynolds numbers has to be carried into the closure and into the mesh.
Taking the usual laminar estimates, the hydrodynamic entry length $0.05\,\mathrm{Re}_{ch}D_h$ is
about 4.5~mm at $\mathrm{Re}_{ch}=50$, under 4\% of the 118~mm channel, so the velocity field is
developed over almost the whole passage. The thermal entry length
$0.05\,\mathrm{Re}_{ch}\mathrm{Pr}\,D_h$ is about 307~mm at the same condition with FC-40 at
$\mathrm{Pr}=67.5$, which is more than twice the channel length. The Graetz number
$\mathrm{Gz}=D_h\mathrm{Re}_{ch}\mathrm{Pr}/L$ is about 52 for the base geometry and about 750 for
the reference configuration. The fin channels are therefore thermally developing over their entire
length in every configuration and at every level of the ladder, and a fully developed Nusselt
constant has no place in the framework's Nusselt closure. It also fixes what the
solver has to get right first, which is why the developing-flow benchmark of Lee and
Garimella~\cite{lee2006} is the lowest rung of the verification ladder.

\subsection{Discretisation, solver and convergence}
\label{sec:discretisation}

The equations are discretised by the finite volume method on a steady, segregated, pressure-based,
double-precision formulation. Pressure and velocity are coupled by the SIMPLE algorithm, which is
the convention of the field by a wide margin: of the twenty numerical studies surveyed here, twelve
name SIMPLE and not one names PISO or SIMPLEC. A coupled solver appears three times, in the anchor,
in B. Zhang et al.\ and in Kim et al.~\cite{chun2026,zhang2026,kim2025}, and is held as the
fallback if SIMPLE proves
slow to converge at the low flow rates, where the momentum equations are weakly coupled and
convergence is dominated by the pressure correction.

Convective terms in the momentum and energy equations are discretised second-order upwind, and
gradients by a least-squares cell-based reconstruction. That choice needs a word, because the
literature is split on it: five of the surveyed studies use second-order
upwind~\cite{dong2026,khoshvaghtaliabadi2025,huang2023ijhmt,cai2019,shahsavar2021} and four use
first-order~\cite{khoshvaghtaliabadi2026oblique,khoshvaghtaliabadi2026component,ni2024,lee2006},
with two sister papers from one group disagreeing with each other. First order is inadmissible
here. The quantity being predicted is a split of mass between two parallel passages whose flow
directions differ by close to ninety degrees at the heat-sink inlet plane, and numerical diffusion
on a mesh that is not aligned with either stream is exactly the error that would smear that split.
Pressure is interpolated by a staggered-mesh scheme of the PRESTO! type, following the anchor,
which selected it for resolving pressure through a heavily blocked region~\cite{chun2026}. Both
sources apply gravity as a body force in the momentum equation, and the pressure interpolation must
be consistent with it.

The numerical framework is implemented in OpenFOAM (version 2406, using \texttt{chtMultiRegionFoam} and \texttt{chtMultiRegionSimpleFoam} for conjugate multi-region solves), parallelised using domain decomposition across 32 cores via hierarchical Scotch partitioning. Pressure-velocity coupling uses the PIMPLE/SIMPLE algorithms, solved with geometric-algebraic multigrid (GAMG) for pressure and preconditioned bi-conjugate gradient (PBiCGStab/DICPCG) solvers for momentum and enthalpy.

Under-relaxation factors for steady conjugate iterations are set to 0.3 on pressure, 0.7 on momentum, and 0.95 on fluid and solid energy equations, consistent with the surveyed literature practice~\cite{dong2026,huang2023ijhmt,khoshvaghtaliabadi2025}.

The mesh is unstructured with prism layers on every wetted wall, refined by proximity in the fin
passages, in the lateral bypass gap and around the flow-control blocks, since those three regions
carry the quantity of interest. Grid independence is established on three monitors rather than one:
heat-source temperature and heat-sink pressure drop, following the anchor~\cite{chun2026}, and the
bypass mass fraction $\Phi_{\mathrm{bypass}}$ of Eq.~(\ref{eq:phibypass}), which is added because
it is the reported quantity most sensitive to how well the gap and the near-wall regions are
resolved and is the one no prior study had to converge. The acceptance threshold is a change below
0.5\% in all three between successive refinements, which is stricter than the anchor's 0.21\% on
temperature and 0.12\% on pressure drop only in the sense that it applies to a third quantity as
well (as detailed in Table~\ref{tab:grid_convergence}, the 720,000-cell Fine mesh and 1.32-million-cell Ultra-Fine mesh achieve $\mathrm{GCI}_{21} \le 0.35\%$ across all monitors, with near-wall $y^+ \approx 0.62$, meeting the ASME standard discretisation criteria).

A solution is accepted as converged only when two conditions hold together. Scaled residuals must
fall below $10^{-4}$ for continuity and momentum and below $10^{-6}$ for energy. That sits at the
strict end of the range the surveyed studies use, where continuity and momentum targets run from
$10^{-3}$ to $10^{-4}$ and energy targets from $10^{-5}$ to
$10^{-7}$~\cite{chun2026,dong2026,huang2023ijhmt,ni2024,kim2025,zhang2026}. Independently of the
residuals, four integral monitors must be stationary over the final several hundred iterations: the
area-averaged case temperature at the heat-sink base, the static pressure drop across the heat
sink, the bypass mass fraction $\Phi_{\mathrm{bypass}}$, and the global energy balance, which must
close to within 0.5\% of the applied load. Requiring a physical quantity to be stationary and not
only a residual to be small follows Dong et al.~\cite{dong2026}, who stop only once the net power
consumption reaches zero and the outlet temperature and pressure are stable, and Cai et
al.~\cite{cai2019}, who require no observable change in surface temperature beyond the residual
criterion. The energy balance is included because the base paper reports no residual target at
all~\cite{muneeshwaran2023} and because a conjugate model with an adiabatic outer boundary has an
exact global check available at no cost. Reporting the bypass fraction as a convergence monitor is
specific to this work: it is the framework's first link, and a solution whose temperature field has
settled while its mass split has not would be useless here even though it would pass every criterion
in the surveyed literature.
